\chapter{RUANG VEKTOR TOPOLOGIS}

\section{Himpunan Seimbang dan Himpunan Penyerap}

\begin{defn} Suatu himpunan $S$ di dalam ruang vektor disebut
 \index{seimbang}%
\df{seimbang} (atau \df{melingkar}) jika $\Dsk S \subseteq S$.
\end{defn}

\begin{notn} Misalkan $\Dsk$ menyatakan cakram satuan tertutup
 \index{D@$\Dsk$ (cakram satuan tertutup di~$\C$)}%
pada bidang kompleks $\{z \in \C\colon \abs z \le 1\}$.
\end{notn}

\begin{prop} Misalkan $B$ suatu subhimpunan seimbang dari sebuah ruang vektor. Jika $\alpha$ dan $\beta$ adalah skalar sedemikian sehingga $\abs\alpha \le \abs\beta$,
maka $\alpha B \subseteq \beta B$.
\end{prop}

\begin{cor} Jika $B$ suatu subhimpunan seimbang dari sebuah ruang vektor dan $\abs\lambda =1$, maka $\lambda B = B$.
\end{cor}

\begin{defn} Jika $S$ suatu subhimpunan dari sebuah ruang vektor, maka himpunan $\Dsk S$ adalah
 \index{seimbang!selubung}%
 \index{selubung!seimbang}%
\df{selubung seimbang} dari~$S$.
\end{defn}

\begin{prop} Misalkan $S$ suatu subhimpunan dari ruang vektor~$V$.  Maka selubung seimbang dari $S$ adalah irisan semua
subhimpunan seimbang dari $V$ yang memuat~$S$; jadi, himpunan itu adalah subhimpunan seimbang terkecil dari $V$ yang memuat~$S$.
\end{prop}

\begin{defn} Suatu subhimpunan $A$ dari ruang vektor $V$
 \index{menyerap}%
\df{menyerap} suatu subhimpunan $B \subseteq V$ jika terdapat $M > 0$ sedemikian sehingga $B \subseteq tA$ setiap kali $\abs t \ge M$.
Himpunan $A$ bersifat
 \index{menyerap}%
\df{menyerap} (atau \df{radial}) jika himpunan itu menyerap setiap himpunan beranggota satu (secara ekuivalen, setiap himpunan hingga).
\end{defn}

Jadi, secara longgar, suatu himpunan bersifat menyerap jika setiap vektor dalam ruang tersebut termuat dalam suatu hasil dilatasi yang sesuai dari himpunan itu.

\begin{exam} Misalkan $B$ adalah gabungan sumbu-$x$ dan sumbu-$y$ dalam ruang vektor riil $\R^2$.  Maka $B$ seimbang, tetapi tidak konveks.
\end{exam}

\begin{exam} Dalam ruang vektor $\C$, sebuah elips dengan salah satu fokusnya di titik asal merupakan contoh himpunan penyerap yang tidak seimbang.
\end{exam}

\begin{exam} Dalam ruang vektor $\C^2$, himpunan $\Dsk \times \{0\}$ seimbang, tetapi tidak menyerap.
\end{exam}

\begin{prop} Suatu himpunan seimbang $B$ dalam ruang vektor $V$ bersifat menyerap jika dan hanya jika untuk setiap $x \in V$ terdapat skalar
$\alpha \ne 0$ sedemikian sehingga $\alpha x \in B$.
\end{prop}











\section{Filter}

Dalam mempelajari ruang vektor topologis, penggunaan bahasa \emph{filter} memudahkan kita mencirikan sifat-sifat topologis
ruang-ruang tersebut.  Bahasa ini merupakan sarana alternatif bagi jaring, tetapi sebagian besar ekuivalen dengannya, dan terutama berguna dalam situasi (nonmetrik) ketika barisan
tidak lagi banyak membantu.

Secara kasar, sebuah \emph{filter} adalah keluarga tak kosong dari himpunan-himpunan tak kosong yang tertutup terhadap irisan berhingga dan pengambilan superhimpunan.

\begin{defn} Suatu keluarga tak kosong $\sfml F$ dari subhimpunan-subhimpunan tak kosong dari suatu himpunan $S$ disebut
 \index{filter!pada suatu himpunan}%
\df{filter} pada $S$ jika
   \begin{enumerate}[label=(\alph*)]
     \item $A$, $B \in \sfml F$ mengakibatkan $A \cap B \in \sfml F$, dan
     \item jika $A \in \sfml F$ dan $A \subseteq B$, maka $B \in \sfml F$.
   \end{enumerate}
\end{defn}

\begin{exam} Salah satu contoh filter yang sederhana (dan, bagi keperluan kita, terpenting) adalah keluarga semua lingkungan suatu titik $a$ di sebuah
ruang topologis.  Keluarga ini adalah
 \index{lingkungan!filter pada suatu titik}%
 \index{filter!lingkungan}%
\df{filter lingkungan} pada~$a$. Filter ini akan
 \index{N@$\sfml N_a$ (filter lingkungan pada~$a$)}%
dilambangkan dengan $\sfml N_a$.
\end{exam}

\begin{defn} Suatu keluarga tak kosong $\sfml B$ yang terdiri atas subhimpunan-subhimpunan tak kosong dari himpunan $S$ disebut
 \index{basis filter}%
\df{basis filter} pada $S$ jika untuk setiap $A$, $B \in \sfml B$ terdapat $C \in \sfml B$ sedemikian sehingga $C \subseteq A \cap B$.

Kita mengatakan bahwa suatu subkeluarga $\sfml B$ dari filter $\sfml F$ adalah
 \index{basis filter!bagi suatu filter}%
 \index{basis!bagi suatu filter}%
\df{basis filter bagi} $\sfml F$ jika setiap anggota $\sfml F$ memuat suatu anggota~$\sfml B$.  (Penting untuk memeriksa bahwa
setiap subkeluarga seperti itu memang merupakan basis filter.)
\end{defn}

\begin{exam} Setiap filter adalah basis filter.
\end{exam}

\begin{exam} Misalkan $a$ sebuah titik dalam ruang topologis.  Setiap basis lingkungan pada $a$ adalah basis filter bagi filter $\sfml N_a$
yang terdiri atas semua lingkungan~$a$. Secara khusus, keluarga semua lingkungan terbuka dari $a$ adalah basis filter bagi~$\sfml N_a$.
\end{exam}

\begin{exam} Misalkan $\sfml B$ suatu basis filter pada himpunan $S$. Maka keluarga $\sfml F$ yang terdiri atas semua himpunan yang memuat anggota $\sfml B$
adalah filter pada $S$, dan $\sfml B$ adalah basis filter bagi~$\sfml F$.  Filter ini akan kita sebut filter yang
 \index{filter!yang dibangkitkan oleh suatu basis filter}%
\df{dibangkitkan oleh}~$\sfml B$.
\end{exam}

\begin{defn} Suatu basis filter (khususnya, suatu filter) $\sfml B$ pada ruang topologis $X$
 \index{konvergensi!suatu filter}%
 \index{filter!konvergensi suatu}%
 \index{konvergensi!suatu basis filter}%
 \index{basis filter!konvergensi suatu}%
 \index{limit!suatu basis filter}%
\df{konvergen} ke sebuah titik $a \in X$ jika setiap anggota $N$ dari $\sfml N_a$, yaitu filter lingkungan pada $a$, memuat suatu anggota~$\sfml B$.
Dalam hal ini, kita
 \index{<arrowto@$\sfml B \sto a$ (konvergensi suatu basis filter)}%
menuliskan $\sfml B \sto a$.
\end{defn}

\begin{prop} Suatu basis filter pada ruang topologis Hausdorff konvergen ke paling banyak satu titik.
\end{prop}

\begin{exam} Misalkan $(x_\lambda)$ sebuah jaring dalam himpunan~$S$ dan $\sfml F$ keluarga semua $A \subseteq S$ sedemikian sehingga jaring $(x_\lambda)$
pada akhirnya berada di~$A$.  Maka $\sfml F$ adalah filter pada~$S$.  Kita akan menyebutnya filter yang
 \index{filter!yang dibangkitkan oleh suatu jaring}%
\df{dibangkitkan oleh} jaring~$(x_\lambda)$.
\end{exam}

\begin{prop} Misalkan $(x_\lambda)$ suatu jaring dalam ruang topologis~$X$, misalkan $\sfml F$ filter yang dibangkitkan oleh~$(x_\lambda)$,
dan $a \in X$.  Maka $\sfml F \sto a$ jika dan hanya jika $x_\lambda \sto a$.
\end{prop}

\begin{exam} Misalkan $\sfml B$ suatu basis filter pada himpunan~$S$.  Misalkan $D$ himpunan semua pasangan terurut $(b,B)$ sedemikian sehingga $b \in B \in \sfml B$.
Untuk $(b_1,B_1)$, $(b_2,B_2) \in D$, definisikan $(b_1,B_1) \le (b_2,B_2)$ jika $B_2 \subseteq B_1$.  Ini membuat $D$ menjadi himpunan terarah.  Sekarang bentuklah
sebuah jaring dalam $S$ dengan mendefinisikan
   \[ x\colon D \sto S\colon (b,B) \mapsto b \]
yakni, tetapkan $x_\lambda = b$ untuk setiap $\lambda = (b,B) \in D$.  Maka $\bigl(x_\lambda\bigr)_{\lambda \in D}$ adalah sebuah jaring dalam~$S$.  Jaring ini adalah jaring yang
 \index{jaring!yang dibangkitkan oleh suatu basis filter}%
 \index{jaring!yang didasarkan pada suatu basis filter}%
\df{dibangkitkan oleh} (atau \df{didasarkan pada}) basis filter~$\sfml B$.
\end{exam}

\begin{prop} Misalkan $\sfml B$ suatu basis filter pada ruang topologis $X$, misalkan $(x_\lambda)$ adalah jaring yang dibangkitkan oleh $\sfml B$, dan misalkan $a \in X$.
Maka $x_\lambda \sto a$ jika dan hanya jika $\sfml B \sto a$.
\end{prop}













